\documentclass[10pt,letterpaper]{article}
\usepackage[margin=0.79in]{geometry}
\usepackage[T1]{fontenc}
\usepackage{lmodern,amsmath,booktabs,enumitem,microtype}
\setlength{\parindent}{0pt}\setlength{\parskip}{4pt}
\setlist{nosep,leftmargin=1.4em}
\title{\vspace{-1.5em}CS 6158 Midterm Progress Report\\[3pt]\large Budget-Constrained Selective Runtime Monitoring}
\author{Shiyun Cheng \and Gengkun Shao}
\date{October 19, 2026}
\begin{document}\maketitle\vspace{-1.7em}
\section{Research Goal and Progress}
Our project asks how to allocate a fixed budget among runtime monitor-creation locations using historical execution frequency, cost, useful traces, violation reports, and redundancy. The primary research question is whether budget-aware selection preserves more violations than simple heuristics or \emph{Valg}, an adaptive online monitor-creation policy, at comparable \emph{realized} overhead. We originally proposed evaluation on held-out executions, including unique-trace and violation recall, runtime and memory, with software evolution and exploration as extensions.

At midterm, we have \textbf{implemented five selection baselines and a reproducible evaluation pipeline}, and performed \textbf{offline selection experiments on historical TraceMOP monitor-profile data for two real Java projects}. Separately, we have a preliminary Valg reproduction dashboard comparing ValgJ and JavaMOP on four Java projects, with a partial ValgT/TraceMOP comparison. These are \emph{distinct experiments}: our budgeted policies have not yet been deployed in the Valg runtime or compared against Valg at matched actual cost.

\section{Implemented Modules and Methods}
The Python selection pipeline reads per-group historical information and produces selected group IDs, budget usage, recall metrics, CSV tables, and plots. Five strategies are implemented: \textbf{random} feasible selection (30 draws per budget in the real-trace evaluation); \textbf{frequency ranking}; \textbf{standalone utility/cost ranking}; \textbf{redundancy-aware greedy}, which maximizes incremental distinct-trace gain per unit cost; and \textbf{additive knapsack}, which maximizes sum of standalone trace counts using a discretized dynamic program. The knapsack surrogate can double-count traces shared by several groups; it is not an exact optimum for distinct-union coverage.

For historical trace IDs $T_i$ attached to a candidate group $i$, a selected set $S$, and a proxy cost $c_i$, the current offline experiment solves the approximation
\begin{equation}
\max_{S}\quad \left|\bigcup_{i\in S}T_i\right|\quad \text{s.t.}\quad \sum_{i\in S}c_i\le B,
\qquad \mathrm{TraceRecall}(S)=\frac{|\cup_{i\in S}T_i|}{|\cup_iT_i|}.
\end{equation}
Here $c_i$ is the workbook's \texttt{cost\_proxy\_event\_checks}, \emph{not execution time}. Candidate rows represent inferred location-by-specification groups; they are not yet verified as independently switchable monitor-creation sites. \textbf{Real-trace violation weighting is not supported yet} because the workbook lacks reliable location-attributed violation sets.

\section{Experiments and Preliminary Results}
\textbf{Historical real-trace offline evaluation.} We processed the partner-provided aggregation for \texttt{javadbf} (153 location-by-specification groups, 278 observed distinct trace IDs) and \texttt{jlemmagen} (66 groups, 2,115 trace IDs). Budgets were 10, 20, 30, 50, and 80\% of summed event-check proxies. Random selection used 30 draws per budget. Table~\ref{tab:results} shows 20\% results; the denominator is the union of trace IDs \emph{already present in historical records}.
\begin{table}[h]\centering\small
\caption{Offline unique-trace recall at 20\% event-check proxy budget (not actual runtime).}\label{tab:results}
\begin{tabular}{lrr}\toprule
Strategy & \texttt{javadbf} & \texttt{jlemmagen} \\\midrule
Random (30-draw mean) & 86.51\% & 65.56\% \\
Frequency & 44.96\% & 88.84\% \\
Utility/cost & 97.84\% & 99.67\% \\
Redundancy-aware greedy & 97.84\% & 99.67\% \\
Additive knapsack & 89.57\% & 99.57\% \\\bottomrule
\end{tabular}\end{table}
The cost-aware approaches preserved most of the previously observed distinct traces under this proxy budget. Greedy and static utility/cost tied on both reported workloads, so these results \emph{do not establish} a benefit from explicitly modeling redundancy. Knapsack can fare worse because its additive objective rewards overlapping traces repeatedly. Frequency ranking is unstable across budgets: static-order packing may select an expensive early group at a larger budget and thereby exclude numerous smaller groups, making independently computed selections non-nested and recall non-monotone. These outcomes are exploratory and depend on the historical aggregation.

\textbf{Valg baseline work.} The partner's reproduction dashboard covers \texttt{javadbf}, \texttt{jlemmagen}, \texttt{route53-infima}, and \texttt{xeger} using no-RV, ValgJ and JavaMOP configurations. The comparison provides a concrete environment for future realized-cost baselines. A trace-oriented ValgT/TraceMOP paired result is available for \texttt{jlemmagen}, but coverage and timing measurements are incomplete across other projects. These baseline runs \emph{cannot yet be directly compared} with offline proxy-budget choices.

\section{Challenges, Risks, and What We Learned}
First, \textbf{violation attribution is missing}: historical reports identify violations, but do not show whether each proposed location would independently detect them. Second, \textbf{proxy cost is not realized overhead}: global instrumentation, shared monitoring work, and interactions between selected groups can invalidate an additive runtime interpretation. Third, the inferred creation location and specification group may not match a controllable creation-site switch; this needs inspection of the monitoring implementation. Fourth, all current real-trace data are from historical execution logs, without a held-out evaluation split. We cannot yet claim generalization, measured memory savings, or superiority over Valg.

\section{Schedule and Deliverables}
We remain on track for the implemented selection-pipeline deliverable; a fully realized-cost-matched Valg comparison remains the main risk. \textbf{Oct. 19--25:} validate group IDs and extraction semantics; reproduce real-trace offline plots and curate workloads. \textbf{Oct. 26--Nov. 1:} implement or test enable/disable of selected monitor-creation sites and collect per-site violation associations. \textbf{Nov. 2--8:} profile representative configurations with repeated runtime and memory measurements. \textbf{Nov. 9--15:} compare budgeted policies and Valg at similar measured overhead using held-out executions. \textbf{Nov. 16--end of term:} robustness analyses, possible revision-transfer/exploration study, final repository and report.

We propose that Shiyun maintain the selection/evaluation pipeline and Gengkun coordinate Valg profiling and instrumentation; the division requires team confirmation. \textbf{Fallback:} if switchable monitors or incremental timing measurements are infeasible, we will report a carefully labeled offline proxy-cost benchmark, explore alternative selective instrumentation, and explicitly narrow claims rather than treating proxy budgets as runtime savings.
\end{document}
